% ============================================================
\section{Problemas encontrados e soluções}\label{sec:problemas}
% ============================================================

Cada problema é registrado com o \textbf{sintoma} (o que se observou), a \textbf{causa}, a
\textbf{solução} adotada e o \textbf{impacto} nos resultados. O objetivo é que quem refizer o
experimento não tropece nos mesmos pontos e que fique claro quais escolhas do método nasceram de
um problema concreto. As entradas abaixo são da fase de planejamento e das primeiras etapas.

\newcommand{\problema}[5]{%
    \subsection{#1}\label{#2}
    \begin{description}[leftmargin=1.2em, itemsep=1pt]
        \item[Sintoma.] #3
        #4
        \item[Impacto.] #5
    \end{description}}

\problema{Limitação de taxa da API da Wikipédia (HTTP 429)}{sec:prob-429}
{No estudo de viabilidade, depois de algumas centenas de requisições seguidas, a API passou a
responder HTTP 429 (\emph{Too Many Requests}) e o download parava no meio.}
{\item[Causa.] A política da Wikimedia limita a taxa de requisições por cliente, sobretudo de
clientes sem um \cod{User-Agent} identificado; buscar artigos um a um multiplica o número de
requisições.
\item[Solução.] O cliente da etapa \etapa{corpus} (\cod{wiki.py}) pede o \emph{wikitext} em
lotes de 50 páginas por requisição (o máximo da API), espera 2,5\,s entre requisições, envia um
\cod{User-Agent} com contato (configurável), usa \cod{maxlag=5}, respeita o cabeçalho
\cod{Retry-After} e, na falta dele, recua exponencialmente (30\,s, 60\,s, \ldots, até 10
minutos). Toda resposta crua fica num cache em disco, então um download interrompido recomeça
de onde parou e a limpeza pode ser refeita sem rede. Os artigos são fixados por \cod{revid} no
\cod{manifest.csv}.}
{Nenhum nos resultados. O download completo fica mais lento (horas, em segundo plano), mas é
retomável e reprodutível.}

\problema{Páginas de usuário e categorias de manutenção na busca por categorias}
{sec:prob-categorias}
{As primeiras buscas nas árvores de categorias traziam páginas de usuário, predefinições,
listas e páginas como ``Artigos sem fontes'' e ``Esboços sobre física''.}
{\item[Causa.] A árvore de categorias da Wikipédia mistura categorias de conteúdo com
categorias de manutenção (artigos que precisam de revisão, esboços, páginas de discussão) e de
organização interna, que também são subcategorias das raízes temáticas.
\item[Solução.] Só entram membros do domínio principal (artigos, \cod{ns=0}); subcategorias cujo
título contém ``Artigos'', ``Esboço'', ``Páginas'', ``Predefinições'', ``Wikipédia'',
``Usuário'', ``Listas'' ou ``Categoria:!'' não são seguidas; títulos ``Lista de\ldots'' são
descartados; a profundidade é limitada a 2 (Seção~\ref{sec:corpus}).}
{Os temas ficam mais limpos. A lista de padrões é heurística e pode deixar passar alguma
categoria de manutenção com outro nome; o descarte de páginas multitema e de biografias
reduz o que sobra.}

\problema{Fragmentação do português no tokenizer do Qwen3}{sec:prob-fragmentacao}
{No estudo de viabilidade, o Qwen3 produziu 1,78 tokens por palavra, e só 28\% dos vértices
eram palavras inteiras (54\% no Tucano, 72\% no GPT-2 português). Quatro das 16 palavras-alvo
candidatas (\emph{corrente}, \emph{onda}, \emph{célula}, \emph{núcleo}) não são um token só.}
{\item[Causa.] O vocabulário BPE do Qwen3 (151.669 ids) é multilíngue e dominado por inglês,
chinês e código; palavras portuguesas menos frequentes, principalmente com acentos, são
divididas em pedaços.
\item[Solução.] As palavras-alvo ficaram restritas às que são um token só (12), contadas só no id
minúsculo precedido de espaço. As categorias de token (Seção~\ref{sec:categorias}) entram em
todas as análises, e as interpretações semânticas se concentram em palavras inteiras de
conteúdo. O custo foi pesado na escolha do modelo (Seção~\ref{sec:mod-escolha}).}
{Menos palavras-alvo e uma fração grande de vértices que são fragmentos. É também um resultado
em si: a P2 compara fragmentos e palavras inteiras.}

\problema{\texttt{hidden\_states[36]} normalizado e \texttt{torch\_dtype} obsoleto no
transformers 5}{sec:prob-hidden}
{Ao planejar a extração, a documentação e o código do transformers 5.16 mostraram que o último
elemento de \cod{hidden\_states} não é a saída bruta do bloco 36. Além disso, carregar o modelo
com \cod{torch\_dtype=} emite um aviso de argumento obsoleto.}
{\item[Causa.] Nos modelos decodificadores do transformers 5, a tupla \cod{hidden\_states}
recebe como último elemento o estado \emph{depois} da RMSNorm final (o que a cabeça de
linguagem lê), no lugar da saída bruta do último bloco; os demais elementos são saídas brutas.
O argumento de tipo passou a se chamar \cod{dtype}.
\item[Solução.] As saídas brutas são capturadas por \emph{hooks} em \cod{layers[0]},
\cod{layers[17]} e \cod{layers[35]} (e nos 36 blocos para o \emph{memmap}), e a versão
normalizada por um \emph{hook} em \cod{norm}, que vira a variante de robustez \cod{L36n}. A
verificação \cod{--verify} confere \cod{norm(hook\_36) == hidden\_states[36]} no modelo real
(Seção~\ref{sec:mod-verificacao}). O carregamento usa \cod{dtype=torch.bfloat16}.}
{Sem os \emph{hooks}, o bloco 36 seria comparado aos blocos 1 e 18 num espaço diferente
(normalizado e reescalado dimensão a dimensão), e a transição $18\to36$ misturaria o efeito do
bloco com o da normalização.}

\problema{Limite por ordem de chegada e artigo igual a parágrafo na simulação}
{sec:prob-simulacao}
{Na simulação do desenho híbrido (\cod{sim\_hybrid.py}), as ocorrências mantidas dos tipos
frequentes (``de'', ``,'') vinham todas dos primeiros parágrafos do núcleo, e cada parágrafo do
núcleo vinha de um artigo diferente.}
{\item[Causa.] O limite $f_{\max}$ era aplicado durante a varredura: a partir da 50.ª ocorrência
de um tipo, as seguintes eram ignoradas. E o núcleo escolhia um parágrafo por artigo.
\item[Solução.] No desenho final (Seção~\ref{sec:amostra}), o núcleo é montado primeiro e o
limite é aplicado depois, por sorteio uniforme entre todas as ocorrências de cada tipo acima do
limite. O núcleo usa 5 artigos $\times$ 2 parágrafos por tema. Um teste verifica que o limite é
aplicado por sorteio.}
{Com o limite por ordem de chegada, os tipos frequentes ficariam concentrados em poucos
parágrafos, o que enviesaria a P3 (tipo e parágrafo ficariam confundidos). Com um parágrafo por
artigo, os rótulos ``artigo'' e ``parágrafo'' da P3 seriam praticamente iguais e não poderiam ser
distinguidos.}

\problema{Conteúdo de referências no texto limpo (\texttt{mwparserfromhell} 0.7.2)}
{sec:prob-ref}
{No corpus do estudo de viabilidade, limpo só com \cod{strip\_code}, apareciam no meio dos
parágrafos títulos de livros, nomes de autores, URLs e datas de acesso, colados à palavra
anterior (``A energiaSilva, Livro, 2001 é\ldots''), e frases com buracos no lugar de fórmulas.}
{\item[Causa.] Na versão 0.7.2 do \cod{mwparserfromhell}, o \cod{strip\_code} remove a marca
\cod{<ref>}, mas mantém o \emph{conteúdo} dela como texto; já o conteúdo de \cod{<math>} some
sem deixar nada no lugar (``a energia é dada por , onde\ldots'').
\item[Solução.] Antes da conversão, a árvore sintática é reescrita: marcas \cod{<ref>} e outras
de bloco são removidas com o conteúdo, e linhas com \cod{<math>} e afins são descartadas
(Seção~\ref{sec:corpus}). Testes com trechos reais de \emph{wikitext} cobrem esses casos.}
{Os números do estudo de viabilidade vêm do corpus com esse defeito; são aproximados, mas o
defeito não muda as conclusões do estudo (ordens de grandeza de parágrafos, tokens e
repetições). O corpus final não tem o problema.}

\problema{\emph{Virtualenv} copiado com o \texttt{pytest} de outro projeto}{sec:prob-venv}
{\cod{.venv/bin/pytest} falhava com erros de importação ou rodava com o interpretador errado,
embora \cod{.venv/bin/python -m pytest} funcionasse.}
{\item[Causa.] O ambiente virtual foi copiado de outro projeto. Os \emph{scripts} de entrada
(\cod{pytest}, \cod{ruff}\ldots) têm na primeira linha (\emph{shebang}) o caminho absoluto do
interpretador, que continuava apontando para o \cod{.venv} do projeto original.
\item[Solução.] Rodar as ferramentas pelo interpretador do projeto
(\cod{.venv/bin/python -m pytest}) ou pelo \cod{uv run}, que usa o ambiente certo; recriar o
ambiente com \cod{uv sync --dev} resolve de vez.}
{Nenhum nos resultados. Fica o alerta de sempre conferir qual interpretador executa os testes.}
